% =============================================================================
\section{La máquina de estados}
% =============================================================================

\begin{frame}{Los cuatro estados y sus transiciones}
  \centering
  \begin{tikzpicture}[font=\footnotesize, >={Stealth[length=2mm]},
      e/.style={draw, rounded corners=3pt, minimum height=1.25cm, text width=2.7cm,
                align=center, line width=0.9pt}]
    \node[e, draw=gris, fill=gris!8] (b) {\textbf{BUSCAR}\\{\scriptsize mira, gira un paso, repite}};
    \node[e, draw=azul, fill=azul!8, right=2.3cm of b] (a) {\textbf{APROXIMAR}\\{\scriptsize avanza centrando el oso}};
    \node[e, draw=verde, fill=verde!8, right=2.3cm of a] (f) {\textbf{ENCONTRADO}\\{\scriptsize detenido}};
    \node[e, draw=naranja, fill=naranja!10, below=1.5cm of a] (h) {\textbf{HUIR}\\{\scriptsize retrocede, gira, avanza}};
    \draw[->, line width=0.9pt] ([yshift=4pt]b.east) -- node[above]{oso confirmado} ([yshift=4pt]a.west);
    \draw[->, line width=0.9pt] ([yshift=-4pt]a.west) -- node[below]{oso perdido} ([yshift=-4pt]b.east);
    \draw[->, line width=0.9pt] (a) -- node[above]{llegó (área)} (f);
    \draw[->, naranja, line width=0.9pt] (b.south) |- node[pos=0.25, left]{mochila} (h.west);
    \draw[->, naranja, line width=0.9pt] ([xshift=10pt]a.south) -- node[right]{mochila} ([xshift=10pt]h.north);
    \draw[->, naranja, line width=0.9pt] (f.south) |- node[pos=0.25, right]{mochila} (h.east);
    \draw[->, gris, dashed, line width=0.9pt] ([xshift=-28pt]h.north) -- ++(0,0.75) -| node[pos=0.22, above]{escape completo} ([xshift=14pt]b.south);
  \end{tikzpicture}

  \vspace{0.5em}
  \relax{\footnotesize
  \begin{columns}[T,onlytextwidth]
    \begin{column}{0.485\textwidth}
      \begin{block}{Lo que hay que mirar}
        Las flechas naranjas salen de \textbf{los tres} estados: la mochila
        interrumpe cualquier cosa. Y \cod{ENCONTRADO} no tiene salida propia:
        el robot queda detenido hasta que termina la corrida o aparece la
        mochila.
      \end{block}
    \end{column}
    \begin{column}{0.485\textwidth}
      \begin{block}{``Confirmado'' no es ``visto''}
        Las transiciones usan las clases \textbf{confirmadas} (3 frames
        seguidos). Lo ``visto'' en un solo frame no cambia el estado: a lo
        sumo hace que el robot se quede quieto para confirmar.
      \end{block}
    \end{column}
  \end{columns}}
\end{frame}

\begin{frame}{Las cuatro reglas que gobiernan todo}
  \begin{block}{1. Huir tiene prioridad incondicional}
    {\small Si la amenaza está confirmada, se huye: aunque el oso esté a la vista,
    aunque falten diez centímetros para llegar, aunque el robot ya esté en
    \cod{ENCONTRADO}.}
  \end{block}
  \begin{block}{2. Una huida no se reevalúa a mitad de camino}
    {\small Las fases 2 y 3 giran y avanzan \textbf{a ciegas}: la cámara pierde la
    mochila y eso es parte del diseño. Si se reevaluara, verla de nuevo
    reiniciaría la maniobra para siempre.}
  \end{block}
  \begin{block}{3. El buzzer es un evento, no un estado}
    {\small La orden sale sólo en el ciclo de la transición que la provoca. Si no, el
    Pi mandaría una orden de buzzer por frame.}
  \end{block}
  \begin{block}{4. Después de huir hay un período sordo a la amenaza}
    {\small No estaba en el diseño: lo encontraron las verificaciones. Al terminar de
    huir la mochila ya no está en cámara, pero el confirmador la sostiene 5
    frames más, y sin los 2 s de sordera esa inercia dispara una segunda huida
    idéntica.}
  \end{block}
\end{frame}

\begin{frame}{\cod{BUSCAR}: mirar, girar un paso, repetir}
  \centering
  \begin{tikzpicture}[font=\scriptsize, x=0.62cm, y=0.5cm]
    % un ciclo: 20 pasos de 0,85 s y la reubicación
    \foreach \i in {0,...,7} {
      \fill[azul!25] (\i*1.7, 0) rectangle (\i*1.7+1.2, 1);
      \fill[naranja!70] (\i*1.7+1.2, 0) rectangle (\i*1.7+1.7, 1);
    }
    \node at (14.3, 0.5) {\dots};
    \fill[azul!25] (15, 0) rectangle (16.2, 1);
    \fill[naranja!70] (16.2, 0) rectangle (16.7, 1);
    \fill[verde!55] (16.7, 0) rectangle (18.7, 1);
    \draw[gris] (0,0) rectangle (18.7,1);
    % leyenda, arriba
    \fill[azul!25] (0.0,1.45) rectangle (0.5,1.95);    \node[right] at (0.55,1.7) {mira quieto 0,6 s};
    \fill[naranja!70] (5.0,1.45) rectangle (5.5,1.95); \node[right] at (5.55,1.7) {gira 0,25 s};
    \fill[verde!55] (9.0,1.45) rectangle (9.5,1.95);   \node[right] at (9.55,1.7) {avanza 1 s};
    \draw[decorate, decoration={brace, mirror, amplitude=4pt}] (0,-0.2) -- node[below=4pt]{barrido: 17 s, unos 20 pasos de $\sim$18\grad{} (una vuelta)} (16.7,-0.2);
    \draw[decorate, decoration={brace, mirror, amplitude=4pt}] (16.7,-0.2) -- node[below=4pt]{reubicación} (18.7,-0.2);
  \end{tikzpicture}

  \vspace{0.3em}
  \relax{\footnotesize
  \begin{columns}[T,onlytextwidth]
    \begin{column}{0.485\textwidth}
      \begin{block}{Las tres partes}
        \begin{itemize}
          \item \textbf{Mirando} (\cod{0, 0}): quieto, con la imagen nítida.
                Unos 6 frames por pausa.
          \item \textbf{Barriendo} (\cod{0, 35}): un pulso de giro de 0,25 s.
          \item \textbf{Reubicando} (\cod{30, 0}): al terminar la vuelta,
                avanza $\sim$18 cm y empieza otra. Sin esto barrería para
                siempre la misma escena.
        \end{itemize}
      \end{block}
    \end{column}
    \begin{column}{0.485\textwidth}
      \begin{block}{Dos atajos}
        \begin{itemize}
          \item \textbf{Frena al ver}: si en un frame aparece el oso o la
                mochila, aunque sin confirmar, \textbf{no gira}: se queda
                quieto para que los frames que confirman salgan nítidos.
          \item \textbf{Busca por donde lo vio}: si viene de perder el oso,
                gira hacia el lado donde estaba (\cod{sentido_barrido}).
        \end{itemize}
      \end{block}
      \begin{block}{Salidas}
        Oso confirmado \hacia{} \cod{APROXIMAR}. Mochila confirmada \hacia{}
        \cod{HUIR}.
      \end{block}
    \end{column}
  \end{columns}}
\end{frame}

\begin{frame}{\cod{BUSCAR}: por qué terminó siendo así}
  \relax{\footnotesize
  \begin{tabular}{L{2.2cm}L{4.6cm}L{6.9cm}}
    \toprule
    \textbf{Cuándo} & \textbf{Cómo era} & \textbf{Qué pasó en el piso} \\
    \midrule
    Diseño original & Girar de corrido a velocidad baja &
      A 124\grad/s la imagen sale \textbf{movida}: YOLO no reconoció el oso en
      ningún frame, aunque le pasó por delante varias veces. \\[2pt]
    30 de septiembre & A pasos: girar 0,4 s, mirar 0,6 s &
      Cada paso renovaba el 61 \% de la imagen ($\sim$32\grad). El oso tiene
      que entrar entero: el margen era mínimo, y el barrido de 9 s no cerraba
      la vuelta. \\[2pt]
    1 de octubre & Pasos de 0,25 s; barrido de 17 s &
      35 \% de la imagen por paso ($\sim$18\grad): el oso queda entero en dos
      pausas seguidas. Encontró el oso a 90\grad{} en 15 a 17 pasos. \\[2pt]
    1 de octubre & Primero mirar, después girar &
      Con el orden inverso, lo primero que hacía el robot era apartar la vista
      de lo que tenía enfrente. Con la mochila adelante, no huía. \\[2pt]
    1 de octubre & Frenar al ver también la mochila &
      La veía dos frames, giraba igual, el tercero salía movido y nunca la
      confirmaba. \\
    \bottomrule
  \end{tabular}}

  \vspace{0.5em}
  \begin{exampleblock}{La medición que ordenó la discusión}
    {\footnotesize \cod{scripts_pi/piso/paso_giro.py} da pulsos de giro y mide con la
    propia cámara cuánto se corrió la escena entre una pausa y la siguiente. Lo
    que decide si el barrido tiene huecos es \textbf{qué fracción de la imagen
    se renueva por paso}, y eso no necesita saber el campo visual.}
  \end{exampleblock}
\end{frame}

\begin{frame}{\cod{APROXIMAR} y \cod{ENCONTRADO}}
  \relax{\footnotesize
  \begin{columns}[T,onlytextwidth]
    \begin{column}{0.54\textwidth}
      \begin{block}{Qué hace \cod{APROXIMAR} en cada ciclo, en este orden}
        \begin{enumerate}
          \item ¿El oso dejó de estar confirmado? \hacia{} anota de qué lado
                se lo vio por última vez y vuelve a \cod{BUSCAR}.
          \item ¿Está confirmado pero no hay caja en \emph{este} frame? \hacia{}
                \textbf{frena} y espera el próximo. Seguir sería andar a
                ciegas.
          \item ¿El área de la caja alcanzó \cod{AREA_OBJETIVO}, \textbf{o} el
                ultrasónico lee 20 cm o menos y lo de adelante es el oso?
                \hacia{} \cod{ENCONTRADO}, con el pitido de hallazgo.
          \item ¿El ultrasónico ve algo a 20 cm o menos que \textbf{no} es el
                oso? \hacia{} desvío, girando hacia el lado contrario.
          \item Si no: \cod{control.perseguir()} \hacia{} avanzar y girar
                hacia el oso.
        \end{enumerate}
      \end{block}
    \end{column}
    \begin{column}{0.43\textwidth}
      \begin{block}{\cod{ENCONTRADO}}
        Manda \cod{0, 0} en cada ciclo. Es terminal para la demo: no vuelve a
        buscar. Lo único que lo saca es la mochila.
      \end{block}
      \begin{alertblock}{Dónde decide y dónde queda}
        \cod{AREA_OBJETIVO} es el tamaño con el que el robot \textbf{decide}
        frenar. Entre la decisión y la detención sigue avanzando (frame viejo
        más inercia). Con 0,35 decide a unos 38 cm y queda a $\sim$30.
      \end{alertblock}
      \begin{block}{Cuánto tarda}
        Desde 1 m: unos 2,5 s de aproximación. Desde que lo confirma de
        frente hasta frenar, 3 s.
      \end{block}
    \end{column}
  \end{columns}}
\end{frame}

\begin{frame}{\cod{HUIR}: tres fases con un solo reloj}
  \centering
  \begin{tikzpicture}[font=\scriptsize, x=2.2cm, y=0.5cm]
    \fill[naranja!30] (0,0) rectangle (1.5,1);
    \fill[naranja!55] (1.5,0) rectangle (2.7,1);
    \fill[naranja!80] (2.7,0) rectangle (5.2,1);
    \fill[gris!25] (5.2,0) rectangle (6.0,1);
    \draw[gris] (0,0) rectangle (6.0,1);
    \node at (0.75,0.5) {1. retroceso};
    \node at (2.1,0.5) {2. giro};
    \node at (3.95,0.5) {3. avance};
    \node at (5.6,0.5) {sordo};
    \foreach \t/\x in {0/0, 1{,}5/1.5, 2{,}7/2.7, 5{,}2/5.2} \node[below] at (\x,0) {\t{} s};
  \end{tikzpicture}

  \vspace{0.2em}
  \relax{\footnotesize
  \begin{tabular}{L{2.0cm}L{1.6cm}L{2.6cm}L{7.4cm}}
    \toprule
    \textbf{Fase} & \textbf{Dura} & \textbf{Comando} & \textbf{Qué hace} \\
    \midrule
    1. Retroceso & 1,5 s & \cod{-40}, $\mp$\cod{25} &
      Marcha atrás $\sim$30 cm, girando hacia el lado \textbf{opuesto} a la
      mochila. Va más lento que el avance: es a ciegas. \\[1pt]
    2. Giro & 1,2 s & \cod{0}, $\mp$\cod{60} &
      Sigue girando \textbf{para el mismo lado} que la fase 1, hasta quedar
      de espaldas. \\[1pt]
    3. Avance & 2,5 s & \cod{60}, \cod{0} &
      Se aleja $\sim$65 cm. Con algo a 30 cm o menos adelante, \textbf{gira en
      el lugar} en vez de avanzar (\cod{0}, \cod{60}). \\[1pt]
    Después & 2,0 s & (en \cod{BUSCAR}) &
      Período sordo: ignora la mochila. \\
    \bottomrule
  \end{tabular}}

  \vspace{0.3em}
  \relax{\footnotesize
  \begin{columns}[T,onlytextwidth]
    \begin{column}{0.485\textwidth}
      \begin{block}{De qué lado gira}
        Al entrar a \cod{HUIR} guarda el signo de \cod{error_x} de la mochila
        (\cod{lado_amenaza}). Si en ese frame no hay caja, usa el último lado
        conocido.
      \end{block}
    \end{column}
    \begin{column}{0.485\textwidth}
      \begin{alertblock}{Lo que sigue a ojo}
        El giro depende del lado: en marcha atrás el robot rota más rápido
        hacia la izquierda que hacia la derecha. Con 1,3 s se pasó de media
        vuelta; 1,2 s no se midió.
      \end{alertblock}
    \end{column}
  \end{columns}}
\end{frame}

\begin{frame}[fragile]{El orden en que decide \texttt{Maquina.paso()}}
\begin{columns}[T,onlytextwidth]
\begin{column}{0.56\textwidth}
\footnotesize
\begin{verbatim}
paso(confirmadas, detecciones, esp, ahora):

  si veo el oso en este frame:
      recuerdo su error_x          # para buscarlo

  si la mochila esta confirmada
     y no estoy huyendo
     y no estoy en el periodo sordo:
        recuerdo de que lado esta
        paso a HUIR (alarma)

  segun el estado:
      BUSCAR      -> _buscar()
      APROXIMAR   -> _aproximar()
      HUIR        -> _huir()
      ENCONTRADO  -> detenido
\end{verbatim}
\end{column}
\begin{column}{0.42\textwidth}
\footnotesize
\begin{block}{El arbitraje, en una tabla}
\begin{tabular}{L{3.25cm}L{1.9cm}}
  \toprule
  \textbf{Situación} & \textbf{Resultado} \\
  \midrule
  Oso y mochila a la vez & Huye \\
  Mochila durante la huida & La ignora \\
  Mochila, 2 s tras huir & La ignora \\
  Mochila, robot detenido & Huye \\
  Oso visto 1 o 2 frames & Queda quieto \\
  Mochila vista 1 o 2 frames & Queda quieto \\
  \bottomrule
\end{tabular}
\end{block}
\begin{block}{Por qué la mochila va primero}
  La comprobación de la amenaza está \textbf{antes} del reparto por estado:
  ningún estado puede ``olvidarse'' de mirarla.
\end{block}
\end{column}
\end{columns}
\end{frame}

\begin{frame}{El ultrasónico en la máquina de estados (desde el 2 de octubre)}
  \relax{\footnotesize
  \begin{tabular}{L{2.6cm}L{2.0cm}L{4.0cm}L{4.9cm}}
    \toprule
    \textbf{Dónde} & \textbf{Umbral} & \textbf{Qué hace} & \textbf{Motivo en pantalla} \\
    \midrule
    \cod{BUSCAR}, pulso de avance & 30 cm &
      Gira lo que dura el pulso, en vez de avanzar &
      \texttt{reubicando: obstáculo (...), giro} \\[1pt]
    \cod{APROXIMAR}, es el oso & 20 cm &
      Da por alcanzado el objetivo &
      \texttt{llegué (sonar)} \\[1pt]
    \cod{APROXIMAR}, no es el oso & 20 cm &
      Gira a tope hacia el lado contrario del oso &
      \texttt{obstáculo ajeno (...), desvío} \\[1pt]
    \cod{HUIR}, fase 3 & 30 cm &
      Gira en el lugar hasta tener el camino libre &
      \texttt{fase 3: obstáculo, esquivo} \\[1pt]
    \cod{HUIR}, fases 1 y 2 & --- &
      Nada: el sensor mira adelante y el robot retrocede o gira &
      --- \\
    \bottomrule
  \end{tabular}}

  \relax{\footnotesize
  \begin{columns}[T,onlytextwidth]
    \begin{column}{0.485\textwidth}
      \begin{block}{Tres reglas}
        \begin{itemize}
          \item \textbf{Dos umbrales}: 20 cm para llegar al oso
                (\cod{DIST_PARADA_CM}), 30 para esquivar
                (\cod{DIST_OBSTACULO_CM}).
          \item \textbf{Retención de 0,3 s}: una lectura ``libre'' suelta no
                alcanza para volver a avanzar.
          \item \textbf{Sin lectura no es distancia cero}: con el sensor
                desconectado, nada de esto actúa.
        \end{itemize}
      \end{block}
    \end{column}
    \begin{column}{0.485\textwidth}
      \begin{block}{¿Es el oso lo que tengo adelante?}
        Sí, si la cámara lo ve con $|error_x| < 0{,}7$ y ya ocupa el 60 \% del
        área objetivo. En la práctica el robot frena antes por tamaño (con el
        sensor en $\sim$28 cm), así que el sonar queda de respaldo.
      \end{block}
      \begin{alertblock}{De dónde salió cada número}
        De tres corridas del 2 de octubre, una por ajuste. Está contado en la
        sección ``El 2 de octubre''.
      \end{alertblock}
    \end{column}
  \end{columns}}
\end{frame}
